\subsection{Scalability Experiments}

To evaluate how the encoding choice affects the Knowledge Graph to quantum
workflow as graph size increases, we extend the running example with two
complementary experiment groups. First, synthetic RDF datasets are used for a
controlled scalability analysis. Second, real Knowledge Graph files from the
professor's QuantumEncoDeco resources are used as a realism check. The purpose
of these experiments is not to claim quantum advantage and not to benchmark an
RDF store or DBpedia as a database system. Instead, the goal is to isolate the
cost of the encoding pipeline itself: RDF parsing, deterministic identifier
mapping, state preparation, circuit construction, local simulation, measurement,
decoding, and resource metadata extraction.

The synthetic datasets use deterministic Turtle triples of the form
\texttt{ex:e0 ex:p0 ex:e1 .}. The controlled sizes are 100, 1,000, 5,000, and
10,000 triples, with predefined numbers of entities and predicates. This
construction makes it possible to compare basis, amplitude, and phase encodings
under a predictable growth pattern. In practice, the default local experiment
uses 100, 1,000, and 5,000 triples, while the 10,000-triple case remains
available as an explicitly requested guarded run when local memory and runtime
permit it.

The real KG files are stored locally in the repository under
\texttt{data/real\_kgs/}. The files used in the realism check are
\texttt{exampleV3.ttl}, \texttt{productsSmall.rdf},
\texttt{DecodedOntologies\_V2.ttl}, and \texttt{Aristotle.xml}. These inputs are
not treated as a benchmark of a specific source collection; rather, they test
whether the same encoding workflow behaves consistently on less controlled RDF
inputs. The comparison remains focused on the encoding families themselves:
basis encoding, amplitude encoding, and phase encoding.

For each dataset and encoding, the experiment records run-level measurements for
RDF parse time, ID mapping time, state preparation time, circuit construction
time, simulation time, measurement time, total runtime, qubit count, and circuit
resource metadata where those quantities are available. Circuit depth and gate
count are reported at multiple levels. The logical metrics are computed on the
original high-level Qiskit circuit, while decomposed and transpiled metrics are
computed where feasible using Qiskit's decomposition and transpilation tools.
This distinction is important because high-level operations such as
\texttt{initialize} and \texttt{unitary} may appear as a single logical
operation even though they should not be interpreted as single hardware-level
gates. Each configuration is repeated several times, using five repetitions by
default, and the raw results are aggregated separately for synthetic datasets,
real KG datasets, and the combined set.

The three encodings expose different trade-offs. Basis encoding represents
subjects, predicates, and objects directly in computational-basis registers, so
its qubit requirement depends on the number of unique entities and predicates.
Amplitude encoding represents triples over padded triple-index states, which can
reduce qubit requirements but still incurs state-preparation and interpretation
costs. Phase encoding also operates over triple indices and applies phase shifts
before an interference step. For synthetic datasets the default marked predicate
is \texttt{http://example.org/p0}, which is guaranteed to occur by construction.
For real KG inputs, the marked predicate is selected from the data, typically as
the most common predicate, so that the real KG phase experiment does not
silently become a no-op. The number and fraction of marked triples are recorded
with each run. In the present implementation, the phase oracle is a dense
diagonal unitary, which is suitable for small local experiments but can become
impractical as the index space grows.

Skipped, timeout, parser-failure, metric-limit, and simulator-limit rows are
reported honestly rather than removed from the results. These rows are important
because they reveal practical bottlenecks in the local KG-to-quantum pipeline.
For example, if a basis-encoded dataset exceeds the configured dense statevector
limit, the experiment still reports the available parsing, mapping, and encoding
metadata. Similarly, if decomposed or transpiled circuit metrics exceed the
configured metric guardrail, the corresponding fields are left unavailable
rather than replaced with artificial values. Such results should be interpreted
as evidence about the feasibility of this local dense-simulation workflow and
the cost of the encoding choices, not as a claim about quantum advantage or all
possible quantum implementations.
